\documentclass[10pt]{article}

% Specify the margins and text size.
\setlength{\textwidth}{6.4in}
\setlength{\textheight}{9.5in}
\setlength{\oddsidemargin}{0pt}
\setlength{\evensidemargin}{0pt}
\setlength{\topmargin}{0pt}
\setlength{\hoffset}{.05in}
\setlength{\voffset}{-1in}

\setlength{\parskip}{5pt}
\setlength{\parindent}{0pt}

% Load some fonts and symbol packages
\usepackage{latexsym}
\usepackage{pifont}       % contains 'star' symbol for counterinsurgency handbook title
\usepackage{yfonts} 
\usepackage{amsmath}
\usepackage{amsfonts}

\usepackage{graphicx}     % actually, this is loaded by pstricks
\usepackage[T1]{fontenc}
\usepackage{ifthen}
\usepackage{pstricks,pst-grad,pst-text,pst-node,multido,pst-plot,calc,pst-3dplot}
%\usepackage[all]{xy}
%\usepackage{animate}

% The hyperref package inserts links into the pdf.
\definecolor{MyLinkColor}{rgb}{.1,.2,1}
\definecolor{MyCiteColor}{rgb}{.1,1,.2}
\definecolor{MyURLColor}{rgb}{.4,.4,.4}
\usepackage[backref=true,pagebackref=false,hyperindex,colorlinks=true,
  linkcolor=MyLinkColor,urlcolor=MyURLColor]{hyperref}


% The tweaklist package is something I found on the web.  It provides a simple interface
% for making changes to spacing used in the itemize and enumerate environments.  Comment
% this out if you don't care to use tweaklist.
\usepackage{tweaklist}
\renewcommand{\itemhook}{\setlength{\parskip}{2pt}\setlength{\parsep}%
{1pt}\setlength{\topsep}{0pt}\setlength{\itemsep}{0pt}}

\newcommand{\U}{\underline{\hspace{5pt}}}

\usepackage{listings}
\newcommand{\Z}{\hphantom{0}}
\newcommand{\HH}{\hspace{20pt} \hphantom{a)} }

\begin{document}
\pagestyle{empty}
\lstset{language=R, showspaces=false, showstringspaces=false}

\href{http://www.shepherd.edu}{\includegraphics[height=1.75cm]{logo-high-res.eps}}
\vspace{-1.79cm}

{\small
\begin{tabular}{cl}
& Math 314\\
& Statistics\\
\hspace{5.28in} & 30 November 2011
\end{tabular}
}
\setlength{\baselineskip}{1.05\baselineskip}

\begin{center}
\textbf{\large  Quiz 8}
\end{center}

1. The spinner for a game-of-chance is tested to see if it is balanced
by spinning it $50$ times. 
Use the data in the table below to test the null hypothesis that that spinner is
fair (i.e.,~that the differences between the observed and expected frequencies are
due to chance). Compute the $P$-value and discuss its significance.

\begin{center}
\begin{tabular}{ccc}
\begin{pspicture}(-2,-2)(2,2)
\pswedge*[linecolor=lightgray](0,0){2}{0}{72}
\pswedge[hatchcolor=lightgray,fillstyle=crosshatch](0,0){2}{72}{144}
\pswedge*[linecolor=white](0,0){2}{144}{216}
\pswedge[hatchcolor=white,fillstyle=crosshatch*,fillcolor=lightgray](0,0){2}{216}{288}
\pswedge[hatchcolor=lightgray,fillstyle=crosshatch](0,0){2}{288}{360}
\pscircle(0,0){2}
\SpecialCoor
\rput{126}(1.6;36){\textbf{A}}
\rput{198}(1.6;108){\textbf{B}}
\rput{270}(1.6;180){\textbf{C}}
\rput{-18}(1.6;252){\textbf{D}}
\rput{54}(1.6;-36){\textbf{E}}
\psset{linewidth=3pt}
\psline{*->}(1.0;-20)(1.4;160)
\pscircle[fillstyle=solid,linewidth=0.04,fillcolor=white](0,0){0.15}
\end{pspicture}
&
\hspace{1in}
&
\begin{pspicture}(-2,-2)(2,2)
\rput(0,0){\begin{tabular}{|c|c|}\hline
        & observed\vphantom{\Large Y}\\
  value & frequency\\[2pt]\hline
    A   &  \Z8\vphantom{\Large Y}\\
    B   &  12\\
    C   & \Z7\\
    D   & \Z9\\
    E   & 14\\\hline
\end{tabular}}
\end{pspicture}
\end{tabular}
\end{center}
\vfill
{\small
\begin{tabular}{|c|rrrrrrr|}\hline
df & 0.5 & 1.0 & 1.5 & 2.0\rput(-19pt,20pt){Table entries are area to the right of the given $\chi^2$ value} & 2.5 & 3.0 & 3.5\vphantom{\Large Y}\\\hline
1 &  0.48 & 0.32 & 0.22 & 0.16 & 0.11 & 0.08 & 0.06\vphantom{\Large Y}\\
2 &  0.78 & 0.61 & 0.47 & 0.37 & 0.29 & 0.22 & 0.17\\
3 &  0.92 & 0.80 & 0.68 & 0.57 & 0.48 & 0.39 & 0.32\\
4  & 0.97 & 0.91 & 0.83 & 0.74 & 0.64 & 0.56 & 0.48\\
5  & 0.99 & 0.96 & 0.91 & 0.85 & 0.78 & 0.70 & 0.62\\
6  & 1.00 & 0.99 & 0.96 & 0.92 & 0.87 & 0.81 & 0.74\\\hline
\end{tabular}}

\eject
{\ }

2. (Hypothetical.)  A county has two high schools, East and West.  
The table below shows the results of a simple random sample of 100 graduates
who attended  these schools then majored in either mathematics or computer science in 
college.  Is  college major independent of the high school attended?
Compute the $P$-value and discuss its significance.

\begin{center}
\begin{tabular}{|l|cc|}\hline
     & East H.S. & West H.S.\vphantom{\Large Y}\\\hline
Math & 38 & 20\vphantom{\Large Y}\\
C.S. & 22 & 20 \\\hline
\end{tabular}
\end{center}
\vfill
\eject

\end{document}

